\begin{example}
    Sea $a \neq 0$, sea $\tau \in \mathbb{R}$ y sea $\sigma > 0$ tal que $\sigma \notin \mathbb{Z}$.
    Entonces, la sucesión $(an^\sigma \log^\tau n)_{n\geq2}$ es equidistribuida módulo $1$.
\end{example}

\begin{resolve}
    Sea $x_n = an^\sigma \log^\tau n$ para cada $n \geq 2$. Definimos $g(x) = a(x+1)^\sigma \log^\tau (x+1)$
    en~$[1,\infty)$, entonces $g(n) = x_{n+1}$ para cada $n \geq 1$. 

    Sea $k = \floor{\sigma}+1$, la derivada $k$-ésima de $g(x)$ es
    \[
        g^{(k)}(x) = a (x+1)^{\sigma-k} \log^{\tau-k}(x+1) \left(P(\log (x+1))\right)
    \]
    donde $P(x)$ es un polinomio de grado $k$ con coeficientes que dependen de $\sigma$ y $\tau$.

    Además, $\sigma - k < 0$ y como $\sigma - k + 1 = \sigma - \floor{\sigma} > 0$ se tiene que
    \[
        \lim_{x \to \infty} g^{(k)}(x) = 0
        \qquad \text{y} \qquad
        \lim_{x \to \infty} x \left|g^{(k)}(x)\right| = \infty
    \]

    Para la monotonia calculamos $g^{(k+1)}$
    \[
        g^{(k+1)}(x) = a (x+1)^{\sigma-k-1} \log^{\tau-k-1}(x+1) \left(Q(\log (x+1))\right)
    \]
    donde $Q(x)$ es un polinomio de grado $k+1$ con coeficientes que dependen de $\sigma$ y $\tau$.

    Debido a que $(x+1)^{\sigma-k-1} \log^{\tau-k-1}(x+1)$ es positivo para cada $x\geq1$ y que
    \[
        \lim_{x \to \infty} Q(x) = -\infty
    \]
    porque el coeficiente del término de mayor grado es $\sigma (\sigma-1) \cdots (\sigma-k) < 0$ 
    entonces existe $x_0 \geq 1$ tal que $Q(\log (x+1)) < 0$ para cada $x \geq x_0$. 
    En consecuencia, $g^{(k+1)}(x)$ es de signo constante para cada $x \geq x_0$, y por tanto, 
    $g^{(k)}(x)$ es monótona para cada $x \geq x_0$.

    Por el teorema \ref{teo:fejer_cap4}, $(g(n))$ es equidistribuida módulo $1$ y 
    por lo tanto, $(x_n)$ también lo es.
\end{resolve}